% L4-4 (2026-11-02): new subsubsection for §6.5 of numerical_results.tex (label subsubsec:num_L4_4), figure environment included.
% Figure: plots/L4-4/aniso_island.png -> figures/L4_4_aniso_island.png (scripts/export_paper_assets.jl L4-4).
% Data: data/exp_pro/L4-4/{lenses,references,y_sensitivity,checks,mesh_diag}.csv (59/70 checks; failures discussed in the text).
% Matrix row L4.4 re-frozen 2026-11-02 08:50 (h = lambda_p/4, nz = 7, frame 2 lambda_p, islands (3,2)/(2,3)/(2.25,1.5) m).
% Open point: A12 (omega = 12, diverging) FE vs PE 9.4 %; the mesh diagnostic shows it is mostly not discretisation.
{\color{red}
\subsubsection{Anisotropic resonator island}\label{subsubsec:num_L4_4}

\paragraph{Test description.} A plate with $C_\rho=0.49~\si{\meter\tothe{5}\per\second\squared}$, $m_\rho=0.01~\si{\meter}$ and $\nu=0$ on water of depth $H=10~\si{\meter}$ carries an elliptic island of resonators ($M_\rho=0.01~\si{\meter\cubed}$ each, $\omega_r=10~\si{\radian\per\second}$), represented by the homogenised resonator field with the density $n(\xvec)=n_0\varphi(\xi)$, $\xi=\sqrt{(x/L_x)^2+(y/L_y)^2}$, a Gaussian $\varphi=e^{-\xi^2/2}$ tapered to zero on $2\le\xi\le2.5$. A plane wave travels along $x$. We solve for the scattered field on half of the domain ($y\ge0$, symmetry imposed on the plate by Nitsche terms), with cells of $\lambda_p/4$, seven graded layers in depth and an absorbing frame of $2\lambda_p$ behind and beside the island. The island load is built from the interpolated incident wave, so that it equals the change of the operator applied to the incident wave to $6\times10^{-13}$. The undamped region extends to $y_\mathrm{end}=12.6$--$16.1~\si{\meter}$. We record $T=\eta/\eta_\mathrm{inc}$ on three lines just behind the island, propagate it back to the island centre with the exact bare-plate propagator and fit the phase: $f=-k_0/\Phi''(0)$ and $\mathcal A=-\ln|T(0)|$. A wide-angle split-step parabolic equation through the exact local wavenumber is processed in the same way; its focal length changes by at most $2.9\%$ when the extent of the line data varies between $y_\mathrm{end}$ and $30~\si{\meter}$. Six cases are considered: islands $(L_x,L_y)=(3,2)$ and $(2,3)~\si{\meter}$ with $n_0M_\rho=0.03~\si{\meter}$ at $\omega=8~\si{\radian\per\second}$ (A8, B8), the first one also with $n_0M_\rho=0.01~\si{\meter}$ (A8w); $(2.25,1.5)~\si{\meter}$ at $\omega=12~\si{\radian\per\second}$, above resonance (A12); and both orientations at $\omega=\omega_r$ with $\zeta=0.05$ and $n_0M_\rho=0.01~\si{\meter}$ (A10, B10), where the island absorbs rather than focuses.

\paragraph{Results.} Below resonance the island is a converging lens and above resonance a diverging one (Fig.~\ref{fig:num_L4_4}). The finite element focal lengths are $12.7$, $34.2$ and $34.2~\si{\meter}$ for A8, B8 and A8w, within $2.8$, $4.7$ and $3.5\%$ of the parabolic equation and within $12\%$ of the ray focus; the ray model misses part of the diffraction by these islands, which are about one wavelength wide. Rotating the island by $90^\circ$ lengthens the focus by a factor $2.7$ (ray: $3.0$). At $\omega=\omega_r$ the optical depths are $\mathcal A=1.326$ and $0.892$, within $1.8$ and $0.6\%$ of the parabolic equation, and their ratio $1.49$ follows the aspect-ratio law $L_x/L_y=1.5$ of Eq.~\eqref{eq:aspect_ratio_law}. Above resonance (A12) the finite element focus, $-8.0~\si{\meter}$, is $9\%$ shorter than the parabolic equation; a refinement from $\lambda_p/4$ to $\lambda_p/5$ on a narrower domain changes it by only $1.4\%$, so the difference is not a discretisation error, and the parabolic reference is itself least reliable here (it differs from the ray focus by $36\%$). Seven or eight layers in depth give the same focal length to $2\times10^{-5}$. The absorbing frame reflects $2$--$3\%$ of the scattered wave, which contributes to a spread of $2.5$--$6.3\%$ between the focal lengths measured on the three lines.

\begin{figure}[t]
    \centering
    \includegraphics[width=\linewidth]{figures/L4_4_aniso_island.png}
    \caption{Anisotropic resonator island. (a) $|\eta/\eta_\mathrm{inc}|$ for $(L_x,L_y)=(3,2)~\si{\meter}$, $\omega=8~\si{\radian\per\second}$ (A8); dashed: island support; white lines: measurement lines. (b) Back-propagated phase at $x=0$: markers, finite element; lines, parabolic equation. (c) Focal lengths from the island centre: finite element, parabolic equation, ray and thin-lens values.}
    \label{fig:num_L4_4}
\end{figure}
}
